\chapter{Gaussian 測度の \texorpdfstring{$W_2$}{W2}}
\label{ch:gaussian-measures}

\section{主結果}
\label{sec:gaussian-main-result}

$\mu_i=\mathcal{N}(m_i,\Sigma_i)$（$i=0,1$）を
$\R^n$ 上の Gaussian 測度とする．ここで
$m_i\in\R^n$ は平均ベクトル，
$\Sigma_i\in\R^{n\times n}$ は対称半正定値な共分散行列である．

\begin{theorem}{Gaussian 測度間の $W_2$ 公式}{gs-gaussian-formula}
  \[
    \boxed{
    \Wass_2(\mu_0,\mu_1)^2
    =
    \norm{m_0-m_1}^2
    +\tr(\Sigma_0)+\tr(\Sigma_1)
    -2\tr\!\left[
      \left(
        \Sigma_0^{1/2}\Sigma_1\Sigma_0^{1/2}
      \right)^{1/2}
    \right]}
  \]
  が成り立つ．
\end{theorem}

\begin{remark}{Bures--Wasserstein 距離}{gs-bures}
  共分散行列に依存する項
  \[
    d_{\mathrm{BW}}(\Sigma_0,\Sigma_1)^2
    :=
    \tr(\Sigma_0)+\tr(\Sigma_1)
    -2\tr\!\left[
      \left(
        \Sigma_0^{1/2}\Sigma_1\Sigma_0^{1/2}
      \right)^{1/2}
    \right]
  \]
  は，現在では Bures--Wasserstein 距離と呼ばれる．
  したがって主公式は
  $\Wass_2^2=\norm{m_0-m_1}^2+d_{\mathrm{BW}}^2$ と分解できる．
\end{remark}

\section{平均と共分散への分解}
\label{sec:mean-covariance-split}

\begin{lemma}{平均の寄与の分離}{gs-mean-split}
  $\mu,\nu\in\Pp{2}(\R^n)$ の平均をそれぞれ $m_\mu,m_\nu$ とし，
  平行移動で中心化した測度を
  \[
    \bar{\mu}:=(x\mapsto x-m_\mu)\pushforward\mu,
    \qquad
    \bar{\nu}:=(y\mapsto y-m_\nu)\pushforward\nu
  \]
  とする．このとき
  \[
    \Wass_2(\mu,\nu)^2
    =
    \norm{m_\mu-m_\nu}^2
    +
    \Wass_2(\bar{\mu},\bar{\nu})^2.
  \]
\end{lemma}

\begin{proof}
  任意の結合 $(X,Y)$ に対し
  $\bar X:=X-m_\mu$，$\bar Y:=Y-m_\nu$ とおくと，
  $\mathbb{E}[\bar X-\bar Y]=0$ なので
  \[
    \begin{aligned}
      \mathbb{E}\norm{X-Y}^2
      &=
      \mathbb{E}\norm{
        (\bar X-\bar Y)+(m_\mu-m_\nu)
      }^2 \\
      &=
      \mathbb{E}\norm{\bar X-\bar Y}^2
      +\norm{m_\mu-m_\nu}^2.
    \end{aligned}
  \]
  結合について下限を取ればよい．
\end{proof}

以後，定理~\ref{thm:gs-gaussian-formula}の証明では
$m_0=m_1=0$ と仮定する．

\begin{lemma}{最適結合は Gaussian に選べる}{gs-gaussian-coupling}
  \blockmeta{uses=prop:gs-optimal-coupling}
  中心 Gaussian 測度
  $\mu_0=\mathcal{N}(0,\Sigma_0)$ と
  $\mu_1=\mathcal{N}(0,\Sigma_1)$ の間には，
  Gaussian である最適結合が存在する．
\end{lemma}

\begin{proof}
  命題~\ref{prop:gs-optimal-coupling}により最適結合
  $(X,Y)$ が存在する．その共分散行列を
  \[
    \Gamma
    =
    \begin{pmatrix}
      \Sigma_0 & K\\
      K^\top & \Sigma_1
    \end{pmatrix},
    \qquad
    K:=\mathbb{E}[XY^\top]
  \]
  とする．同じ共分散行列 $\Gamma$ をもつ中心 Gaussian ベクトル
  $(X_{\mathrm{G}},Y_{\mathrm{G}})$ を取ると，
  その周辺分布は $\mu_0,\mu_1$ であり，
  \[
    \mathbb{E}\norm{X_{\mathrm{G}}-Y_{\mathrm{G}}}^2
    =
    \tr(\Sigma_0)+\tr(\Sigma_1)-2\tr(K)
    =
    \mathbb{E}\norm{X-Y}^2.
  \]
  したがって Gaussian 結合も最適である．
\end{proof}

\section{共分散行列の最適化}
\label{sec:covariance-optimization}

補題~\ref{lem:gs-gaussian-coupling}より，中心 Gaussian の場合は
交差共分散 $K$ を選ぶ有限次元問題へ帰着する：
\[
  \Wass_2(\mu_0,\mu_1)^2
  =
  \tr(\Sigma_0)+\tr(\Sigma_1)
  -2\sup_K\tr(K),
\]
ただし
\[
  \begin{pmatrix}
    \Sigma_0 & K\\
    K^\top & \Sigma_1
  \end{pmatrix}
  \succeq0
\]
を満たす $K$ の上で上限を取る．

\begin{lemma}{実行可能な交差共分散}{gs-cross-covariance}
  $\Sigma_0,\Sigma_1$ が正定値なら，
  ブロック行列が半正定値であることと
  \[
    K=\Sigma_0^{1/2}C\Sigma_1^{1/2},
    \qquad
    \norm{C}_{\mathrm{op}}\leq1
  \]
  を満たす行列 $C$ が存在することは同値である．
\end{lemma}

\begin{proof}
  合同変換により
  \[
    \begin{pmatrix}
      \Sigma_0 & K\\
      K^\top & \Sigma_1
    \end{pmatrix}
    \succeq0
    \quad\Longleftrightarrow\quad
    \begin{pmatrix}
      I & C\\
      C^\top & I
    \end{pmatrix}
    \succeq0,
    \qquad
    C:=\Sigma_0^{-1/2}K\Sigma_1^{-1/2}.
  \]
  右辺に Schur 補行列を適用すると
  $I-C^\top C\succeq0$，すなわち
  $\norm{C}_{\mathrm{op}}\leq1$ と同値である．
\end{proof}

\begin{lemma}{トレース最大化}{gs-trace-max}
  任意の実正方行列 $A$ に対し
  \[
    \sup_{\norm{C}_{\mathrm{op}}\leq1}\tr(CA)
    =
    \tr\!\left((A^\top A)^{1/2}\right).
  \]
\end{lemma}

\begin{proof}
  $A=U\Lambda V^\top$ を特異値分解とする．
  $D:=V^\top C U$ とおけば $\norm{D}_{\mathrm{op}}\leq1$ であり，
  \[
    \tr(CA)=\tr(D\Lambda)
    \leq\sum_{i=1}^n\Lambda_{ii}.
  \]
  $C=VU^\top$ とすれば等号を達成する．右辺は
  $\tr((A^\top A)^{1/2})$ に等しい．
\end{proof}

\begin{proof}[定理~\ref{thm:gs-gaussian-formula}の証明]
  まず $\Sigma_0,\Sigma_1$ が正定値で，平均が $0$ とする．
  補題~\ref{lem:gs-cross-covariance}とトレースの巡回性から
  \[
    \begin{aligned}
      \sup_K\tr(K)
      &=
      \sup_{\norm{C}_{\mathrm{op}}\leq1}
      \tr(\Sigma_0^{1/2}C\Sigma_1^{1/2}) \\
      &=
      \sup_{\norm{C}_{\mathrm{op}}\leq1}
      \tr(C\Sigma_1^{1/2}\Sigma_0^{1/2}).
    \end{aligned}
  \]
  補題~\ref{lem:gs-trace-max}を
  $A=\Sigma_1^{1/2}\Sigma_0^{1/2}$ に適用すると
  \[
    \sup_K\tr(K)
    =
    \tr\!\left[
      \left(
        \Sigma_0^{1/2}\Sigma_1\Sigma_0^{1/2}
      \right)^{1/2}
    \right].
  \]
  これを輸送コストの式へ代入すれば，中心 Gaussian に対する公式を得る．
  平均が非零の場合は補題~\ref{lem:gs-mean-split}を使う．

  共分散行列が特異な場合は
  $\Sigma_i+\varepsilon I$ に対する公式を用いる．
  独立な $Z\sim\mathcal{N}(0,I)$ を加える結合から
  \[
    \Wass_2\!\left(
      \mathcal{N}(m_i,\Sigma_i+\varepsilon I),
      \mathcal{N}(m_i,\Sigma_i)
    \right)
    \leq\sqrt{n\varepsilon}
  \]
  なので，$\varepsilon\downarrow0$ とすれば一般の場合が従う．
\end{proof}

\section{最適輸送写像}
\label{sec:gaussian-map}

\begin{proposition}{Gaussian 間の最適輸送写像}{gs-gaussian-map}
  $\Sigma_0$ が正定値なら
  \[
    T(x)
    :=
    m_1+A(x-m_0),
    \qquad
    A
    :=
    \Sigma_0^{-1/2}
    \left(
      \Sigma_0^{1/2}\Sigma_1\Sigma_0^{1/2}
    \right)^{1/2}
    \Sigma_0^{-1/2}
  \]
  は $T\pushforward\mu_0=\mu_1$ を満たし，
  $(\Id,T)\pushforward\mu_0$ は $\Wass_2$ の最適結合である．
\end{proposition}

\begin{proof}
  $A$ は対称半正定値であり
  \[
    A\Sigma_0A^\top=\Sigma_1
  \]
  を満たす．したがって $X\sim\mathcal{N}(m_0,\Sigma_0)$ なら
  $T(X)\sim\mathcal{N}(m_1,\Sigma_1)$ である．
  この結合のコストを展開すると
  定理~\ref{thm:gs-gaussian-formula}の右辺に一致するので最適である．
\end{proof}

\section{特殊な場合}
\label{sec:gaussian-special-cases}

\begin{proposition}{可換な共分散行列}{gs-commuting-covariances}
  $\Sigma_0\Sigma_1=\Sigma_1\Sigma_0$ なら
  \[
    \Wass_2(\mu_0,\mu_1)^2
    =
    \norm{m_0-m_1}^2
    +
    \norm{\Sigma_0^{1/2}-\Sigma_1^{1/2}}_{\mathrm{F}}^2,
  \]
  ここで $\norm{\cdot}_{\mathrm{F}}$ は Frobenius ノルムである．
\end{proposition}

\begin{proof}
  可換な対称行列は同時対角化できるので
  \[
    \left(
      \Sigma_0^{1/2}\Sigma_1\Sigma_0^{1/2}
    \right)^{1/2}
    =
    \Sigma_0^{1/2}\Sigma_1^{1/2}.
  \]
  主公式へ代入すればよい．
\end{proof}

\begin{example}{一次元 Gaussian}{gs-one-dimensional-gaussian}
  $\mu_i=\mathcal{N}(m_i,\sigma_i^2)$ とすると
  \[
    \Wass_2(\mu_0,\mu_1)^2
    =
    (m_0-m_1)^2+(\sigma_0-\sigma_1)^2.
  \]
\end{example}

\begin{example}{二次元の中心 Gaussian}{gs-two-dimensional-gaussian}
  $m_0=m_1=0$ かつ $\Sigma_0,\Sigma_1\in\R^{2\times2}$ とする．
  $2\times2$ 半正定値行列 $B$ に対する恒等式
  $\tr(B^{1/2})^2=\tr(B)+2\sqrt{\det(B)}$ を用いると
  \[
    \Wass_2(\mu_0,\mu_1)^2
    =
    \tr(\Sigma_0)+\tr(\Sigma_1)
    -2\sqrt{
      \tr(\Sigma_0\Sigma_1)
      +2\sqrt{\det(\Sigma_0\Sigma_1)}
    }.
  \]
\end{example}
